% !TeX program = xelatex
\documentclass[11pt,a4paper,fontset=fandol]{ctexart}
\usepackage{guide}
\title{Guide · 样式与命令手册}
\subtitle{把精致的排版，变成简单的日常书写}
\author{LaTeX 模板组件库}
\institute{中文笔记 · 技术手册 · 项目说明}
\date{2026 年 9 月}
\guideheader{Guide · 样式与命令手册}{Template / Tex}
\begin{document}
\maketitle
\begin{infobox}[从 main.tex 开始]
复制 \texttt{main.tex} 与 \texttt{guide.sty} 到新文件夹，修改标题和正文，
使用 XeLaTeX 或 Tectonic 编译。本手册展示各项命令的实际排版效果。
\end{infobox}
\tableofcontents
\vfill
\muted{样式提取自 Xiaolong Luo (Aaron) 的 phd-survival-guide，遵循 MIT 许可。
模板保留原作的主题配色、图标列表与代码样式，并封装常用提示框和中文排版。}

\clearpage
\iconsection{palette}{文字与配色}
\iconsubsection{pen-fancy}{用颜色表达层次}
普通正文用于叙述，\themecolor{主题色用于关键词}，
\important{红色粗体用于重点}，\success{绿色用于完成状态}，
\muted{灰色用于补充信息}。\highlight{黄色底纹}适合短语，
\badge{进行中}和\badge[emeraldgreen]{已完成}适合状态标签。
\begin{texcode}
\themecolor{keyword}       \important{key point}
\success{done}         \muted{additional note}
\highlight{short phrase}
\highlight[blue!10]{custom background}
\badge{in progress}    \badge[emeraldgreen]{done}
\textcolor{royalpurple}{custom text color}
\end{texcode}
\begin{highlightbox}
较长的重点段落放入 \texttt{highlightbox}，正文会根据页面宽度自动换行。
需要标题时，可以选择下一页的提示框。
\end{highlightbox}
\subsection{原作六色}
\begin{center}
\renewcommand{\arraystretch}{1.25}
\begin{tabular}{lll}
\toprule
色样 & 颜色名称 & RGB \\
\midrule
\swatch{darkblue} & \texttt{darkblue} & 0, 0, 139 \\
\swatch{lightgray} & \texttt{lightgray} & 240, 240, 240 \\
\swatch{harvardcrimson} & \texttt{harvardcrimson} & 165, 28, 48 \\
\swatch{emeraldgreen} & \texttt{emeraldgreen} & 0, 155, 119 \\
\swatch{royalpurple} & \texttt{royalpurple} & 102, 51, 153 \\
\swatch{goldenyellow} & \texttt{goldenyellow} & 255, 193, 37 \\
\bottomrule
\end{tabular}
\end{center}
\subsection{调整主题}
导言区写入 \texttt{\textbackslash guidetheme\{royalpurple\}}，即可统一调整
标题、默认图标、强调文字和链接。也可以先用
\texttt{\textbackslash definecolor} 定义自己的颜色，再传给此命令。

\clearpage
\iconsection{icons}{图标与列表}
\emoji{book} 书籍\quad \emoji{lightbulb} 想法\quad
\emoji{rocket} 行动\quad \emoji{flask} 实验\quad
\emoji{chart-line} 结果\quad \emoji{graduation-cap} 学术

\emoji[harvardcrimson]{heart} 关怀\quad \emoji{code} 代码\quad
\emoji{calendar-check} 日程\quad \emoji{users} 合作\quad
\emoji{brain} 思考\quad \emoji{check-circle} 完成
\begin{texcode}
\emoji{lightbulb}
\emoji[harvardcrimson]{heart}
\iconsection{book}{Section title}
\iconsubsection{flask}{Subsection title}
\iconsection*{rocket}{Unnumbered title}
\end{texcode}
\muted{图标使用 Font Awesome 5 名称。\texttt{\textbackslash emoji} 与
\texttt{\textbackslash icon} 等价；它们输出可以着色的矢量字形。
带图标标题会在目录中保留图标，PDF 书签使用标题文字。}
\subsection{七种语义列表}
\begin{itemize}
  \gooditem 已完成：数据整理和实验记录。
  \baditem 待解决：一项尚未通过的验证。
  \toolitem 工具：完成工作所需的资源。
  \ideaitem 想法：值得进一步研究的假设。
  \warningitem 注意：需要及时处理的问题。
  \staritem 重点：最值得记住的结论。
  \arrowitem 下一步：可以马上执行的行动。
\end{itemize}
\begin{texcode}
\begin{itemize}
  \gooditem Completed task.
  \baditem Unresolved issue.
  \toolitem Useful tool.
  \ideaitem A new idea.
  \warningitem Pay attention.
  \staritem Key result.
  \arrowitem Next action.
\end{itemize}
\end{texcode}

\clearpage
\iconsection{comment-dots}{提示框与重点段落}
\begin{infobox}[信息 · infobox]
用来解释背景、定义和阅读方式。
\end{infobox}
\begin{tipbox}[建议 · tipbox]
用来记录可操作的方法、经验和改进建议。
\end{tipbox}
\begin{warningbox}[注意 · warningbox]
用来标注前提、约束和需要读者关注的事项。
\end{warningbox}
\begin{importantbox}[重点 · importantbox]
用来强调关键结论或需要优先处理的问题。
\end{importantbox}
\begin{reflectionbox}[思考 · reflectionbox]
用来整理假设、观察和进一步讨论的问题。
\end{reflectionbox}
\begin{guidebox}[emeraldgreen]{check-circle}{自定义 · guidebox}
自定义颜色、图标和标题；框内可以写段落、列表和公式。
\end{guidebox}
\begin{texcode}
\begin{tipbox}[My suggestion]
Write paragraphs or a list here.
\end{tipbox}

\begin{guidebox}[royalpurple]{brain}{Reflection}
Write your idea here.
\end{guidebox}
\end{texcode}
\muted{五种语义框的标题参数可省略，使用默认中文标题。所有提示框支持跨页。}

\clearpage
\iconsection{code}{代码与表格}
\subsection{Python：行号与语法高亮}
\begin{pythoncode}[caption={数据归一化示例}]
def normalize(values):
    # Keep experiment inputs comparable.
    total = sum(values)
    return [v / total for v in values]

scores = normalize([2, 3, 5])
print(scores)
\end{pythoncode}
\subsection{Bash：浅蓝背景与左侧强调线}
\begin{bashcode}
# Compile the document
tectonic main.tex
\end{bashcode}
\subsection{三个可直接使用的环境}
使用 \texttt{pythoncode}、\texttt{bashcode}、\texttt{texcode} 环境包围代码。
可选参数传给 \texttt{listings}，例如 \texttt{[caption=示例,firstnumber=10]}。
\begin{texcode}
\begin{pythoncode}[caption={Example}]
print("Hello, research!")
\end{pythoncode}

\codefile{analysis.py}
\codefile[bashstyle]{run.sh}
\end{texcode}
\subsection{简洁的三线表}
\begin{center}
\renewcommand{\arraystretch}{1.25}
\begin{tabularx}{\linewidth}{lXr}
\toprule
阶段 & 记录内容 & 状态 \\
\midrule
问题定义 & 明确目标、输入和预期输出 & \success{完成} \\
方案设计 & 记录方法、假设和可验证的判断 & \themecolor{进行中} \\
结果整理 & 汇总证据、结论和下一步行动 & \muted{待开展} \\
\bottomrule
\end{tabularx}
\end{center}
\muted{表格使用 \texttt{booktabs} 和 \texttt{tabularx}；图片使用
\texttt{graphicx}，数学公式使用 \texttt{amsmath}，均已加载。}

\clearpage
\iconsection{swatchbook}{科研配色与快速定制}
\subsection{原作九组色板}
\begin{center}
\renewcommand{\arraystretch}{1.55}
\begin{tabular}{ll}
\toprule
名称 & 预览 \\
\midrule
\texttt{Blues} & \palette{Blues} \\
\texttt{Viridis} & \palette{Viridis} \\
\texttt{Cividis} & \palette{Cividis} \\
\texttt{RdBu} & \palette{RdBu} \\
\texttt{PuOr} & \palette{PuOr} \\
\texttt{BrBG} & \palette{BrBG} \\
\texttt{Set1} & \palette{Set1} \\
\texttt{OkabeIto} & \palette{OkabeIto} \\
\texttt{Dark2} & \palette{Dark2} \\
\bottomrule
\end{tabular}
\end{center}
\begin{texcode}
\palette{Viridis}
\swatch{emeraldgreen}
\hexswatch{0072B2}
\definepalette{Project}{234E70,009E73,E69F00}
\palette{Project}
\end{texcode}
\subsection{新建文档的基本骨架}
\begin{texcode}
\documentclass[11pt,a4paper,fontset=fandol]{ctexart}
\usepackage{guide}
\title{My research notes}
\subtitle{Questions, evidence and actions}
\author{Your name}
\date{\today}
\guideheader{Research notes}{Project name}
\begin{document}
\maketitle
\iconsection{book}{Introduction}
Write your content here.
\end{document}
\end{texcode}
\end{document}
